\section{总结与延伸}

\subsection{核心要点回顾}

本文系统梳理了 Agent Skill 的五种设计模式：

\begin{enumerate}
    \item \textbf{Tool Wrapper}：按需加载，注入特定库知识
    \item \textbf{Generator}：模板驱动，确保输出一致性
    \item \textbf{Reviewer}：checklist 外置，实现可插拔审查
    \item \textbf{Inversion}：先问后做，门控式信息收集
    \item \textbf{Pipeline}：硬性检查点，强制顺序执行
\end{enumerate}

\begin{warningbox}{最重要的启示}
不要试图将复杂而脆弱的指令塞进一个 system prompt 中。将工作流分解，应用合适的结构化模式，才能构建出可靠的 Agent。
\end{warningbox}

\subsection{延伸思考}

\begin{itemize}
    \item \textbf{Agent Skills 规范}是开源的，且已被 ADK 原生支持。开发者可以直接基于这些模式开始构建。
    \item 随着更多 Agent 工具采纳统一的 skill 格式，\textbf{跨工具的 skill 复用}将成为可能。一个设计良好的 skill 可以同时在 Claude Code、Gemini CLI 和 Cursor 中使用。
    \item 模式的组合潜力值得进一步探索：例如 Inversion + Generator + Reviewer 的三层组合，可以构建一个"先收集需求、再生成文档、最后自审"的完整工作流。
\end{itemize}

\subsection{拓展阅读}

\begin{itemize}
    \item Google Agent Development Kit (ADK) 官方文档
    \item Agent Skills 开源规范
    \item 原文作者：@Saboo\_Shubham\_ 和 @lavinigam
\end{itemize}

%% --- Fin del contenido --- %%

\end{document}
